\documentclass{article}
\usepackage{amsmath, amssymb}
\usepackage{geometry}
\geometry{a4paper, margin=1in}

\title{Mathematical Foundations of Spline Geometry Analysis}
\author{}
\date{}

\begin{document}
\maketitle

This document provides the theoretical justification for why we use the geometric metric $G_e$ (based on Deviation-From-Linearity and Curvature) to interpret KAN edge computations, and under what conditions it isolates biological signals.

\section*{Proposition 1: Reference-Induced Deformation Decomposition}

Let the empirical computation along an edge $(i, j)$ be defined as:
\begin{equation}
\Phi_{ij}(x) = W_{ij}^{base}\text{SiLU}(x) + \sum_k w_{ij,k}B_k(x)
\end{equation}

When evaluating the non-linearity of this function, we compute a geometric test statistic $G(f, X)$ (e.g., integrated squared second derivative or deviation from linear fit) over a defined input distribution $X$.

\textbf{The Decomposition Hypothesis:}\\
We assert that the observed geometry can be linearly decomposed into two components:
\begin{equation}
G(f, X) = S(f) + D(X, f)
\end{equation}
Where:
\begin{enumerate}
    \item $S(f)$ is the \textbf{intrinsic structural component}. This represents the true, task-associated mathematical transformation (e.g., thresholding, saturation) required to solve the biological objective.
    \item $D(X, f)$ is the \textbf{reference/input-dependent deformation}. This is the artifactual geometry induced by how the function $f$ interpolates the specific input density $X$, constrained by the flexibility of the B-splines.
\end{enumerate}

\section*{The Role of Null Calibration}

When evaluating $G(f, X)$ across two different cohorts (e.g., HCP $\to$ CHCP), a naive evaluation yields $G(f_{real}, X_{CHCP})$. If this value changes, it is mathematically impossible to distinguish whether the true task mechanism $S(f)$ changed, or merely the input density $X$ shifted, altering $D(X, f)$.

To resolve this identifiability failure, we train an empirical null model $f_{null}$ on label-shuffled data. Because $Y_{null}$ contains no task information, the null model requires no intrinsic biological mechanism:
\begin{equation}
S(f_{null}) \approx 0
\end{equation}

However, $f_{null}$ is trained on the exact same input distribution $X$, utilizing the identical spline flexibility and regularization dynamics. Thus, it fully absorbs the input-dependent deformation:
\begin{equation}
G(f_{null}, X) = D(X, f_{null})
\end{equation}

By constructing the \textbf{Standardized Separation} estimator (or simply evaluating the difference):
\begin{align}
\Psi &= G(f_{real}, X) - G(f_{null}, X) \\
\Psi &= [S(f_{real}) + D(X, f_{real})] - D(X, f_{null})
\end{align}

\textbf{Crucial Theoretical Guarantee:}\\
Assuming the baseline deformation is statistically independent of the true signal mechanism (i.e., $D(X, f_{real}) \approx D(X, f_{null})$ in expectation), the calibration yields:
\begin{equation}
\Psi \approx S(f_{real})
\end{equation}

This proves that the calibrated metric isolates the intrinsic structural behavior relative to the null expectation, removing the confounding effects of $X$. If $\Psi > 0$, the edge exhibits non-linearity exceeding what is forced by the mere routing of the input distribution.

\end{document}
